% ============================================================
%  SCHOLA DOMESTICA — Magister Mathematica
%  Level 1 | Week 14 | Day 3
%  Topic: Fact Families 0–10 — Complete Set
% ============================================================
\documentclass[12pt, letterpaper]{article}
\usepackage{sd_style}

\renewcommand{\lessonweek}{14}
\renewcommand{\lessonnum}{3}
\renewcommand{\lessontitle}{Fact Families 0–10: Complete Set}

% IMAGE PROMPT (Beatrix Potter watercolour style):
% A mother duck with ten ducklings on a quiet pond; she gathers six under her wings
% while four swim in a little cluster nearby. Soft morning mist. White background, no text.

\begin{document}

\lessonheader{Week 14}{3}{Fact Families 0--10: Complete Set}

{\small Name:\ansblanklg \hfill Date:\ansblank}
\goldrule

\begin{instructbox}
Every set of three numbers has a \textbf{fact family} of four sentences.\\
The \textbf{largest number} is always the \emph{whole}; the two smaller numbers are the \emph{parts}.\\[4pt]
Example — Family \{2, 8, 10\}:\quad
$2+8=10$\quad $8+2=10$\quad $10-8=2$\quad $10-2=8$
\end{instructbox}

\bigskip
\noindent{\color{sd-blue}\large\textbf{Part I — Complete the Family}} \hfill \textit{(Problems 1–16)}

\medskip\noindent\textit{The whole is given. Write all four sentences.}

\bigskip
% Family {1,9,10}
\noindent\textbf{Family A: 1, 9, 10}

\noindent\begin{tabular}{@{}p{4cm}p{4cm}p{4cm}p{4cm}@{}}
\prob $1+9=\ansblank$ & \prob $9+1=\ansblank$ &
\prob $10-9=\ansblank$ & \prob $10-1=\ansblank$
\end{tabular}

\medskip
% Family {3,7,10}
\noindent\textbf{Family B: 3, 7, 10}

\noindent\begin{tabular}{@{}p{4cm}p{4cm}p{4cm}p{4cm}@{}}
\prob $3+7=\ansblank$ & \prob $7+3=\ansblank$ &
\prob $10-7=\ansblank$ & \prob $10-3=\ansblank$
\end{tabular}

\medskip
% Family {4,6,10}
\noindent\textbf{Family C: 4, 6, 10}

\noindent\begin{tabular}{@{}p{4cm}p{4cm}p{4cm}p{4cm}@{}}
\prob $4+6=\ansblank$ & \prob $6+4=\ansblank$ &
\prob $10-6=\ansblank$ & \prob $10-4=\ansblank$
\end{tabular}

\medskip
% Family {2,7,9}
\noindent\textbf{Family D: 2, 7, 9}

\noindent\begin{tabular}{@{}p{4cm}p{4cm}p{4cm}p{4cm}@{}}
\prob $2+7=\ansblank$ & \prob $7+2=\ansblank$ &
\prob $9-7=\ansblank$ & \prob $9-2=\ansblank$
\end{tabular}

\bigskip\bigskip
\noindent{\color{sd-blue}\large\textbf{Part II — Missing Number}} \hfill \textit{(Problems 17–22)}

\medskip\noindent\textit{Think of the fact family to find the missing number.}

\bigskip
\noindent\begin{tabular}{@{}p{4.5cm}p{4.5cm}p{4.5cm}@{}}
\prob $6+\ansblank=10$ & \prob $\ansblank+8=10$ & \prob $10-\ansblank=3$ \\[14pt]
\prob $9-\ansblank=7$ & \prob $\ansblank-6=4$ & \prob $\ansblank-9=1$
\end{tabular}

\bigskip
\begin{checkbox}
\textbf{\color{sd-green}Special Case: Doubles!}\\[3pt]
When both parts are the same (like $5+5=10$), the family has only \textbf{two} sentences:\\[3pt]
\hspace{2em}$5+5=10$ \qquad $10-5=5$\\[3pt]
\textit{Can you name the other double families in our range?} \hfill \textit{(Hint: 0+0, 1+1, 2+2, 3+3, 4+4)}
\end{checkbox}

\end{document}
